%% kapa: exam=1 round=31 year=2020 period=2 form=A subject=accounting subject_ko=회계학 seq=33 number=73
%% source: Q-Net 감정평가사 기출문제 (cst00309) artlSeq=5207219 file=extracted/1008/5207219/A형/A형/2020년도 제31회 감정평가사 1차 시험_2교시_A형.pdf
%% conversion: kordoc 4.24.51 -> tools/gen1.py (auto); review=proofread-sonnet; flags=none
\begin{question}{73}
(주)감평은 단일공정을 통해 단일제품을 생산하고 있으며, 선입선출법에 의한 종합원가계산을 적용하고 있다. 직접재료는 공정 초에 전량 투입되고, 가공원가는 공정 전반에 걸쳐 균등하게 발생한다. (주)감평의 20x1년 기초재공품은 10,000단위(가공원가 완성도 40\%), 당기착수량은 30,000단위, 기말재공품은 8,000단위(가공원가 완성도 50\%)이다. 기초재공품의 직접재료원가는 ￦170,000이고, 가공원가는 ￦72,000이며, 당기투입된 직접재료원가와 가공원가는 각각 ￦450,000과 ￦576,000이다. 다음 설명 중 옳은 것은? (단, 공손 및 감손은 발생하지 않는다.)
\begin{choices}
  \item 기말재공품원가는 ￦192,000이다.
  \item 가공원가의 완성품환산량은 28,000단위이다.
  \item 완성품원가는 ￦834,000이다.
  \item 직접재료원가의 완성품환산량은 22,000단위이다.
  \item 직접재료원가와 가공원가에 대한 완성품환산량 단위당원가는 각각 ￦20.7과 ￦20.3이다.
\end{choices}
\end{question}
\answer{1}
